%% kapa: exam=1 round=29 year=2018 period=1 form=A subject=real_estate subject_ko=부동산학원론 seq=27 number=107
%% source: Q-Net 감정평가사 기출문제 (cst00309) artlSeq=5202808 file=extracted/1008/5202808/2018년 제29회 감정평가사 1차시험 기출문제/2018년 제29회 감정평가사 1차시험 기출문제/1교시 A형.pdf
%% conversion: kordoc 4.24.51 -> tools/gen1.py (auto); review=proofread-sonnet; flags=none
\begin{question}{107}
부동산마케팅에 관한 설명으로 옳지 않은 것은?
\begin{choices}
  \item 부동산 공급자가 부동산시장을 점유하기 위한 일련의 활동을 시장점유마케팅전략이라 한다.
  \item AIDA 원리는 소비자가 대상 상품을 구매할 때까지 나타나는 심리 변화의 4단계를 의미한다.
  \item 시장점유마케팅전략에 해당되는 STP전략은 시장세분화(segmentation), 표적시장 선정(targeting), 포지셔닝(positioning)으로 구성된다.
  \item 고객점유마케팅전략에 해당되는 4P MIX전략은 유통경로(place), 제품(product), 위치선점(position), 판매촉진(promotion)으로 구성된다.
  \item 고객점유마케팅전략은 AIDA 원리를 적용하여 소비자의 욕구를 충족시키기 위해 수행된다.
\end{choices}
\end{question}
\answer{4}
